% Generated draft: requires mathematical and source review.
\documentclass[11pt]{article}
\usepackage[margin=25mm]{geometry}
\usepackage{fontspec}
\setmainfont{Latin Modern Roman}
\usepackage{amsmath,amssymb,mathtools,unicode-math}
\setmathfont{Latin Modern Math}
\usepackage{enumitem,needspace,xurl,hyperref}
\setlength{\parindent}{0pt}
\setlength{\parskip}{5pt}
\setlength{\emergencystretch}{3em}
\newcommand{\N}{\mathbb{N}}
\newcommand{\Z}{\mathbb{Z}}
\newcommand{\Q}{\mathbb{Q}}
\newcommand{\R}{\mathbb{R}}
\newcommand{\C}{\mathbb{C}}
\newcommand{\F}{\mathbb{F}}
\newcommand{\A}{\mathbb{A}}
\newcommand{\PP}{\mathbb{P}}
\newcommand{\E}{\mathbb{E}}
\newcommand{\eps}{\varepsilon}
\newcommand{\dd}{\,\mathrm{d}}
\begin{document}
\section*{UnsolvedMath ID: 10400057}
\textbf{Source problem code:} catalogue entry
\subsection{Definitions and mathematical statement}
Let $\F_q$ be a finite field of characteristic $p$. A finite-type knot invariant with values in $\F_q$ is extended to singular knots by $v(K_\times)=v(K_+)-v(K_-)$ in that field, and has degree at most $d$ if it vanishes on knots with more than $d$ transverse double points.

The source asks to construct a Kontsevich, that is universal Vassiliev, invariant with finite-field coefficients. In a standard diagram-valued formulation one seeks
\[
Z_{\F_q}:\mathcal K\longrightarrow\widehat{\mathcal A}_{\F_q}
\]
with a specified degree-completed diagram target, such that every degree-$\leq d$ finite-type invariant factors linearly through its degree-$\leq d$ truncation:
\[
v(K)=\ell_v\bigl(Z_{\F_q,\leq d}(K)\bigr).
\]
The truncation must itself be a finite-type invariant of the required degree, with the prescribed leading-symbol convention. This is not merely an arbitrary encoding of all knot types.

The target diagram space, its coefficient-dependent relations and symbol realization must be fixed; their characteristic-zero universality cannot be imported without proof in characteristic $p$. Rational Kontsevich coefficients cannot in general be reduced modulo $p$ by dividing by denominators that vanish in $\F_q$. A construction valid in a few degrees is also different from one compatible in all degrees.

The record does not specify whether it asks for one finite field, every prescribed field, or a uniform construction; oriented/framed or long/closed knot conventions; the exact diagram relations/normalization; or further multiplicative/structural properties. The display is a conventional universal-invariant target identifying the core factorization requirement, not a claimed existence theorem or an assumed positive-characteristic diagram isomorphism.
\subsection{Short English statement}
Construct an all-degree universal finite-type knot invariant over a specified finite field. Its diagram target and characteristic-$p$ universality must be proved, rather than obtained by undefined reduction of rational denominators.
\subsection{Sources}
\begin{itemize}
\item Local UnsolvedMath registry, permanent ID 10400057, exact source code catalogue entry, Problem 3.7; no proposer/year supplied. Registry status open, classification OPEN-TRIAGE. The registry title is truncated.
\item Ohtsuki, \emph{Problems on invariants of knots and 3-manifolds} (2002), Problem 3.7, PDF page 54, \url{https://msp.org/gtm/2002/04/gtm-2002-04-024s.pdf}. Supplied pdf-font-structure extraction accessed July 29, 2026; provenance distinguishes 231 bold numbered items. No PDF was accessed here.
\item The supplied summary identifies finite-field Kontsevich universality but gives no precise later construction, obstruction or theorem reference; its not-obviously-resolved assessment is not an independent current audit.
\item Provenance flags current primary-literature status and redistribution terms NEEDS\_REVIEW.
\end{itemize}
\end{document}
